\documentclass[11pt]{article}

\usepackage{amsmath,amssymb}
\usepackage{xurl}
\usepackage[hidelinks]{hyperref}
\setlength{\emergencystretch}{2em}

% Shared commands for notes in docs/paper-gaps/.
% Notation follows the LDT paper (references/ldt-paper/).

\newcommand{\Lean}{\textsc{Lean}}
\newcommand{\leanid}[1]{\textsf{#1}}
\newcommand{\ghissue}[1]{\href{https://github.com/LionSR/MIPStarRE/issues/#1}{GitHub issue \##1}}
\newcommand{\ghpr}[1]{\href{https://github.com/LionSR/MIPStarRE/pull/#1}{GitHub PR \##1}}

% --- Paper-compatible notation ---
\newcommand{\F}{\mathbb{F}}
\newcommand{\N}{\mathbb{N}}
\newcommand{\C}{\mathbb{C}}
\newcommand{\tr}{\operatorname{tr}}
\newcommand{\ot}{\otimes}
\newcommand{\E}{\mathop{\bf E\/}}
\newcommand{\bx}{\mathbf{x}}
\newcommand{\by}{\mathbf{y}}
\newcommand{\bu}{\mathbf{u}}
\newcommand{\bv}{\mathbf{v}}

% Approximation relation (paper uses \approx_\eps and \simeq_\zeta)
\newcommand{\approxeps}{\approx_{\varepsilon}}
\newcommand{\simgeq}{\simeq_{\zeta}}

% Polynomial spaces (paper uses \polyfunc, \polysub, \polymeas)
\newcommand{\polyfunc}[3]{\operatorname{Poly}(#1,#2,#3)}
\newcommand{\polysub}[3]{\operatorname{PolySub}(#1,#2,#3)}
\newcommand{\polymeas}[3]{\operatorname{PolyMeas}(#1,#2,#3)}


\title{Naimark dilation: deliberate paper-gap for the full bipartite assembly}
\author{}
\date{2026-05-08}

\begin{document}
\maketitle

\section{At a glance}

\begin{itemize}
  \item \textbf{Difficulty: hard if pursued, low if left as a paper-gap.}  The
  one-measurement Naimark step is straightforward in finite dimensions, but the
  full tensor-product assembly of per-question dilations into a single bipartite
  Naimark statement is left out of this formalization by design.
  \item \textbf{Estimated weight, if formalized:} roughly \textbf{6--10 theorem
  interfaces} and \textbf{600--1200 lines}, dominated by tensor-product
  bookkeeping for ancillary registers and the lifting of states from the small
  to the enlarged space.
  \item \textbf{Mathlib/project split, if formalized:} approximately
  \textbf{30\% Mathlib} and \textbf{70\% project-local}.  Mathlib supplies
  finite-dimensional Hilbert-space tensor products and orthogonal projection
  machinery; the project supplies the question-indexed assembly and the
  bipartite distance-preservation interface.
  \item \textbf{Key Mathlib inputs, if formalized:} finite-dimensional inner
  product spaces, tensor products of finite-dimensional Hilbert spaces, and
  matrix block constructions for projections.
\end{itemize}

\section{Key theorem forms}

The formalization isolates two related but distinct theorem forms.
\begin{enumerate}
  \item \textbf{Source theorem form: Naimark dilation
  (\Cref{thm:naimark-paper}).}  For every collection of bipartite
  submeasurements $\{A^x\}_{x\in Q_A}$, $\{B^y\}_{y\in Q_B}$ on a finite
  Hilbert space~$\mathcal H$, there is an enlarged Hilbert space $\widetilde
  {\mathcal H}\supseteq\mathcal H$, an isometry $V:\mathcal H\to\widetilde
  {\mathcal H}$, and projective measurements $\{\widetilde A^x\}$,
  $\{\widetilde B^y\}$ on $\widetilde{\mathcal H}\otimes\widetilde{\mathcal H}$
  with
  \[
    \langle\psi|A^x_a\otimes B^y_b|\psi\rangle
    =\langle V\psi|\widetilde A^x_a\otimes \widetilde B^y_b|V\psi\rangle.
    \tag{\texttt{thm:naimark}}
  \]
  This is the displayed paper statement at
  \path{references/ldt-paper/orthonormalization.tex}, lines~36--115.
  \item \textbf{Adapter theorem form: per-question one-measurement
  Naimark.}  For each question $x$ and submeasurement $A^x$ there exists a
  one-measurement Naimark dilation on the enlarged space
  $\mathcal H\otimes(\C\oplus\C^{|\textsf{Outcome}|})$ that preserves all
  single-outcome expectations:
  \[
    \tau\bigl(\rho\,A^x_a\bigr)
    =\widetilde\tau\bigl(\widetilde\rho\,\widetilde A^x_{(a,\mathrm{some})}\bigr).
  \]
  This is the form actually formalized; see the Lean structure
  \leanid{NaimarkStatement}.%
  \footnote{Defined at
  \doclink{MIPStarRE/LDT/MakingMeasurementsProjective/Statements.lean}.}
  \item \textbf{Downstream theorem form.}  The downstream uses inside the LDT
  pipeline never call the bipartite tensor product of two Naimark dilations;
  they call the per-question single-outcome marginal-preservation conclusion at
  Alice or Bob separately, with the unmodified state on the unaffected side.
  This is supplied directly by \leanid{leftMarginalPreservation} and
  \leanid{rightMarginalPreservation}.
\end{enumerate}

\section{The source assertion}

We examine \cite{Ji2020LowIndividualDegree}, Theorem~\ref{thm:naimark}.%
\footnote{Tracked alongside ledger~\ghissue{1379}.  The paper passage is
\path{references/ldt-paper/orthonormalization.tex}, lines~36--115, with the
helper \texttt{lem:naimark-helper} used in the proof at lines~117--272 and the
genuine non-preservation example \texttt{ex:easy-but-long} at lines~189--272.}

In the notation of the paper, the source assumes
\[
  \mathcal H=\text{a finite-dimensional Hilbert space},
  \quad \{A^x\}_{x\in Q_A},\,\{B^y\}_{y\in Q_B}\text{ submeasurements on }\mathcal H,
\]
and concludes the bipartite preservation identity above.  The proof proceeds
question by question via \texttt{lem:naimark-helper}; the example
\texttt{ex:easy-but-long} demonstrates that the dilation does \emph{not}
preserve the state-dependent distance $\approx_\delta$, which is why downstream
arguments in the LDT paper carry the dilation only as far as the
single-outcome marginals.

\section{Comparison with the formalization}

The blueprint and formal development together record the per-question
one-measurement dilations and the single-outcome marginal preservation, but
omit the global tensor-product assembly into one bipartite statement.  The
formal development proves the adapter form
\[
  \forall x,a,\rho,\;
  \tau\bigl(\rho A^x_a\bigr)
  =\widetilde\tau\bigl(\widetilde\rho_x\,(\widetilde A^x)_{(a,\mathrm{some})}\bigr),
\]
where $\widetilde\rho_x$ is the one-measurement-lifted density at question $x$.
The corresponding statement for Bob is symmetric.  The bipartite identity
of the source theorem form is \emph{not} proved.

This is sufficient for downstream LDT use: every consumer reduces to the
per-side single-outcome marginal-preservation conclusion, applied either to
Alice's submeasurements with $B$ left untouched or vice versa.  No consumer
requires the simultaneous bipartite dilation; in particular, no consumer
requires the global isometry $V$ acting on the joint state.

\section{Why this is recorded as a paper-gap rather than a producer task}

Three considerations support leaving the bipartite tensor assembly out of the
formalization:
\begin{itemize}
  \item \textbf{Distance non-preservation.}  Example~\texttt{ex:easy-but-long}
  shows that the dilation does not preserve $\approx_\delta$, so any
  bipartite statement formulated in terms of state-dependent distances would
  need to track an explicit error inflation that is not used downstream.
  \item \textbf{Downstream redundancy.}  The downstream uses listed above all
  reduce to one of the per-question marginal-preservation conclusions.  Adding
  the bipartite assembly would not eliminate any obligation in the rest of the
  pipeline.
  \item \textbf{Tensor bookkeeping cost.}  The bipartite assembly is dominated
  by tensor-product bookkeeping (lifting the joint state, transporting density
  matrices through the isometry, transporting submeasurements through the
  enlarged register on each side) which is independent of the LDT-specific
  content.
\end{itemize}

\section{Conclusion}

The Lean structure \leanid{NaimarkStatement} records the per-question
one-measurement Naimark data and the single-outcome marginal-preservation
conclusions used by downstream LDT lemmas.  The full bipartite tensor-product
assembly of the paper's Naimark theorem is intentionally not formalized.  The
classification \texttt{T} in the audit ledger may be promoted to \texttt{E} on
the basis of this note, since the paper origin is now cited at the
\leanid{NaimarkStatement} declaration site and the deliberate gap is recorded
here.

\end{document}
